\documentclass[11pt,a4paper]{article}
\usepackage[margin=2.2cm]{geometry}
\usepackage{amsmath,amssymb,mathtools}
\usepackage[T1]{fontenc}
\usepackage[utf8]{inputenc}
\usepackage{lmodern}
\usepackage{hyperref}
\usepackage{listings}
\lstset{basicstyle=\ttfamily\small,breaklines=true,frame=single}

\title{Derivazione Simbolica Van Vleck delle Costanti Quartiche Watson\\
\large Output in termini del tensore $\tau$ e dei blocchi $H_{mn}$}
\author{Codex + workflow SymPy custom}
\date{\today}

\begin{document}
\maketitle

\section{Obiettivo}
Si vuole derivare simbolicamente il blocco quartico rotazionale dell'Hamiltoniana efficace rovibrazionale,
partendo dalla decomposizione standard in blocchi \(H_{mn}\),
\begin{equation}
H = H_{20} + H_{02} + H_{12} + H_{30} + H_{22} + H_{40},
\end{equation}
con
\begin{align}
H_{20} &= \sum_i \hbar\omega_i\left(a_i^\dagger a_i + \frac12\right),\\
H_{02} &= \text{blocco rotazionale rigido},\\
H_{12} &= \frac12 \sum_{\alpha\beta k} \mu_{\alpha\beta,k} q_k J_\alpha J_\beta,\\
H_{22} &= \frac14 \sum_{\alpha\beta kl} \mu_{\alpha\beta,kl} q_k q_l J_\alpha J_\beta,\\
H_{30} &= \frac16 \sum_{ijk} \Phi_{ijk} q_i q_j q_k,\\
H_{40} &= \frac1{24} \sum_{ijkl} \Phi_{ijkl} q_i q_j q_k q_l.
\end{align}
Le coordinate normali sono
\begin{equation}
q_i = \sqrt{\frac{\hbar}{2\omega_i}}\,(a_i + a_i^\dagger),
\end{equation}
con commutatore bosonico
\begin{equation}
[a_i,a_j^\dagger]=\delta_{ij}.
\end{equation}

\section{Scelta perturbativa standard}
Per rispettare il power-counting standard usato in spettroscopia rovibrazionale:
\begin{align}
 H^{(0)} &= H_{20} + H_{02},\\
 H^{(1)} &= H_{12} + H_{30},\\
 H^{(2)} &= H_{22} + H_{40}.
\end{align}

Con questa scelta, il blocco quartico rotazionale dell'Hamiltoniana efficace si organizza
naturalmente per classi di contributi:
\begin{itemize}
\item \((H_{12},H_{12})\),
\item \((H_{22})\),
\item \((H_{12},H_{30})\),
\item \((H_{30},H_{30})\),
\item \((H_{40})\).
\end{itemize}
Il codice filtra esplicitamente l'algebra simbolica mantenendo solo queste classi, oltre ai
vincoli di parita bosonica e grado rotazionale quartico.

\section{Implementazione}
File implementato:
\begin{center}
\texttt{/Users/vincenzobarone/centrifugal/derive\_watson\_quartic\_vanvleck.py}
\end{center}

\subsection{Algebra operatoriale}
Il codice usa una rappresentazione esplicita di monomi operatoriali:
\begin{itemize}
\item parola vibrazionale (tuple di \texttt{("a",k)} e \texttt{("ad",k)}),
\item parola rotazionale non commutativa (tuple di \(J_x,J_y,J_z\)).
\end{itemize}

Sono implementate:
\begin{itemize}
\item moltiplicazione di espressioni operatoriali,
\item commutatore \([A,B]=AB-BA\),
\item normal ordering bosonico automatico via riscrittura ricorsiva
\(a_i a_j^\dagger \to a_j^\dagger a_i + \delta_{ij}\).
\end{itemize}

\subsection{Van Vleck (Lie transform)}
Si costruisce
\begin{equation}
K = e^{S} H e^{-S} = H + [S,H] + \frac{1}{2!}[S,[S,H]] + \frac{1}{3!}[S,[S,[S,H]]] + \cdots
\end{equation}
con
\begin{equation}
S = \sum_{n\ge1} S^{(n)}.
\end{equation}
A ogni ordine \(m\), \(S^{(m)}\) e scelto imponendo la cancellazione della parte off-diagonal
(vibrazionalmente non conservante) di \(K^{(m)}\).

\subsection{Selezione termini fisici}
Dopo normal ordering si conserva solo:
\begin{itemize}
\item proiezione su ground state vibrazionale: \(\langle 0|\cdots|0\rangle\),
\item blocco quartico rotazionale: monomi con quattro operatori \(J\).
\end{itemize}

Il primo output fisico estratto non sono le costanti Watson, ma i sei coefficienti del tensore
quartico simmetrico \(\tau\),
\begin{equation}
(\tau_{xxxx},\tau_{yyyy},\tau_{zzzz},\tau_{xxyy},\tau_{xxzz},\tau_{yyzz}),
\end{equation}
che moltiplicano rispettivamente
\begin{equation}
J_x^4,\ J_y^4,\ J_z^4,\ J_x^2J_y^2,\ J_x^2J_z^2,\ J_y^2J_z^2.
\end{equation}
Le costanti Watson sono poi ottenute come semplice proiezione lineare del tensore \(\tau\)
sulla riduzione desiderata (A oppure S).

\section{Comandi usati}
\begin{lstlisting}
python3 derive_watson_quartic_vanvleck.py --max-order 2
python3 derive_watson_quartic_vanvleck.py --max-order 3
python3 derive_watson_quartic_vanvleck.py --max-order 4
\end{lstlisting}

\section{Convenzione compatta per le formule}
Per rendere leggibili i risultati simbolici del run usato (1 modo, ansatz diagonale):
\begin{align}
m_x &= \mu1\_0, & m_y &= \mu1\_4, & m_z &= \mu1\_8,\\
n_x &= \mu2\_0, & n_y &= \mu2\_4, & n_z &= \mu2\_8,\\
\phi_3 &= \phi3\_0, & \phi_4 &= \phi4\_0, & \omega &= \omega0.
\end{align}

\section{Risultati}
\subsection{Ordine 2}
\paragraph{Contributo \((H_{12},H_{12})\).}
Nel run simbolico corrente (1 modo, ansatz diagonale), al secondo ordine sopravvive solo la
classe \((H_{12},H_{12})\). I contributi a \(\tau\) sono
\begin{align}
\tau_{xxxx} &= -\frac{m_x^2}{8\omega^2}, &
\tau_{yyyy} &= -\frac{m_y^2}{8\omega^2}, &
\tau_{zzzz} &= -\frac{m_z^2}{8\omega^2},\\
\tau_{xxyy} &= -\frac{m_xm_y}{4\omega^2}, &
\tau_{xxzz} &= -\frac{m_xm_z}{4\omega^2}, &
\tau_{yyzz} &= -\frac{m_ym_z}{4\omega^2}.
\end{align}

Proiettando questo tensore sulla riduzione Watson A si ottiene
\begin{align}
D_J &= \frac{m_x m_y}{8\omega^2},\\
D_{JK} &= \frac{-2m_xm_y + m_xm_z + m_ym_z}{8\omega^2},\\
D_K &= \frac{m_xm_y - m_xm_z - m_ym_z + m_z^2}{8\omega^2},\\
d_1 &= \frac{m_z(m_x-m_y)}{8\omega^2},\\
d_2 &= \frac{m_x^2 - m_y^2 - 2m_z(m_x-m_y)}{16\omega^2}.
\end{align}
Quindi, a questo ordine, le costanti Watson quartiche dipendono esclusivamente dal contributo
\((H_{12},H_{12})\). Non compaiono ancora contributi da \(H_{22}\), \(H_{30}\) o \(H_{40}\).

\subsection{Ordine 3}
Nel run simbolico corrente: nessun nuovo contributo quartico nel ground state
\(\bigl(K^{(3)}\) non aggiunge termini \(J^4\bigr)\). In termini dei blocchi \(H_{mn}\), questo
significa che le classi candidate con numero totale dispari di inserzioni non generano,
nel setup corrente, un nuovo contributo effettivo al tensore \(\tau\).

\subsection{Ordine 4}
\paragraph{Contributi \((H_{22})\), \((H_{12},H_{30})\), \((H_{30},H_{30})\), \((H_{40})\).}
Al quarto ordine compaiono invece i contributi nuovi attesi:
\begin{itemize}
\item termini lineari in \(n_\alpha\), che derivano da \(H_{22}\);
\item termini lineari in \(\phi_4\), che derivano da \(H_{40}\);
\item termini quadratici in \(\phi_3\), che derivano dalla classe \((H_{30},H_{30})\);
\item termini misti \(m_\alpha n_\beta \phi_3\), che riflettono la classe \((H_{12},H_{30})\)
  combinata con i contributi di ordine successivo del Lie transform.
\end{itemize}

Nel report qui sotto si riporta direttamente la proiezione Watson A del risultato quartico:
\begin{align}
D_J &= \frac{-2\hbar m_x m_y\omega^2\phi_4 + 3\hbar m_x m_y\phi_3^2 -3\hbar m_x n_y\omega^2\phi_3 -3\hbar m_y n_x\omega^2\phi_3 + \hbar n_x n_y\omega^4 + 8 m_x m_y\omega^5}{64\omega^7},\\
D_{JK} &= \frac{1}{64\omega^7}\Big(4\hbar m_x m_y\omega^2\phi_4 -6\hbar m_x m_y\phi_3^2 -2\hbar m_x m_z\omega^2\phi_4 +3\hbar m_x m_z\phi_3^2 \\
&\qquad +6\hbar m_x n_y\omega^2\phi_3 -3\hbar m_x n_z\omega^2\phi_3 -2\hbar m_y m_z\omega^2\phi_4 +3\hbar m_y m_z\phi_3^2 \\
&\qquad +6\hbar m_y n_x\omega^2\phi_3 -3\hbar m_y n_z\omega^2\phi_3 -3\hbar m_z n_x\omega^2\phi_3 -3\hbar m_z n_y\omega^2\phi_3 \\
&\qquad -2\hbar n_x n_y\omega^4 + \hbar n_x n_z\omega^4 + \hbar n_y n_z\omega^4 -16 m_x m_y\omega^5 +8 m_x m_z\omega^5 +8 m_y m_z\omega^5\Big),\\
D_K &= \frac{1}{64\omega^7}\Big(-2\hbar m_x m_y\omega^2\phi_4 +3\hbar m_x m_y\phi_3^2 +2\hbar m_x m_z\omega^2\phi_4 -3\hbar m_x m_z\phi_3^2 \\
&\qquad -3\hbar m_x n_y\omega^2\phi_3 +3\hbar m_x n_z\omega^2\phi_3 +2\hbar m_y m_z\omega^2\phi_4 -3\hbar m_y m_z\phi_3^2 \\
&\qquad -3\hbar m_y n_x\omega^2\phi_3 +3\hbar m_y n_z\omega^2\phi_3 -2\hbar m_z^2\omega^2\phi_4 +3\hbar m_z^2\phi_3^2 \\
&\qquad +3\hbar m_z n_x\omega^2\phi_3 +3\hbar m_z n_y\omega^2\phi_3 -6\hbar m_z n_z\omega^2\phi_3 \\
&\qquad +\hbar n_x n_y\omega^4 -\hbar n_x n_z\omega^4 -\hbar n_y n_z\omega^4 +\hbar n_z^2\omega^4 \\
&\qquad +8 m_x m_y\omega^5 -8 m_x m_z\omega^5 -8 m_y m_z\omega^5 +8 m_z^2\omega^5\Big),\\
d_1 &= \frac{1}{64\omega^7}\Big(-2\hbar m_x m_z\omega^2\phi_4 +3\hbar m_x m_z\phi_3^2 -3\hbar m_x n_z\omega^2\phi_3 +2\hbar m_y m_z\omega^2\phi_4 \\
&\qquad -3\hbar m_y m_z\phi_3^2 +3\hbar m_y n_z\omega^2\phi_3 -3\hbar m_z n_x\omega^2\phi_3 +3\hbar m_z n_y\omega^2\phi_3 \\
&\qquad +\hbar n_x n_z\omega^4 -\hbar n_y n_z\omega^4 +8 m_x m_z\omega^5 -8 m_y m_z\omega^5\Big),\\
d_2 &= \frac{1}{128\omega^7}\Big(-2\hbar m_x^2\omega^2\phi_4 +3\hbar m_x^2\phi_3^2 +4\hbar m_x m_z\omega^2\phi_4 -6\hbar m_x m_z\phi_3^2 \\
&\qquad -6\hbar m_x n_x\omega^2\phi_3 +6\hbar m_x n_z\omega^2\phi_3 +2\hbar m_y^2\omega^2\phi_4 -3\hbar m_y^2\phi_3^2 \\
&\qquad -4\hbar m_y m_z\omega^2\phi_4 +6\hbar m_y m_z\phi_3^2 +6\hbar m_y n_y\omega^2\phi_3 -6\hbar m_y n_z\omega^2\phi_3 \\
&\qquad +6\hbar m_z n_x\omega^2\phi_3 -6\hbar m_z n_y\omega^2\phi_3 +\hbar n_x^2\omega^4 -2\hbar n_x n_z\omega^4 -\hbar n_y^2\omega^4 +2\hbar n_y n_z\omega^4 \\
&\qquad +8 m_x^2\omega^5 -16 m_x m_z\omega^5 -8 m_y^2\omega^5 +16 m_y m_z\omega^5\Big).
\end{align}

Quindi al quarto ordine la proiezione Watson contiene esattamente i contributi attesi dai blocchi
\(H_{22}\), \(H_{40}\), \((H_{30},H_{30})\) e \((H_{12},H_{30})\), oltre al termine base
\((H_{12},H_{12})\).

\section{Nota su generalita e simmetrie}
I risultati sopra sono ottenuti in un setup simbolico ridotto (1 modo, componenti diagonali)
per mantenere formule chiuse e leggibili. L'algoritmo implementativo e pero generale:
\begin{itemize}
\item supporta operatori non commutativi rotazionali,
\item impone normal ordering bosonico esatto,
\item separa diagonal/off-diagonal per costruire \(S^{(n)}\),
\item puo essere esteso a piu modi e tensori completi.
\end{itemize}

Per imporre esplicitamente le simmetrie complete
\(\Phi_{ijk}=\Phi_{\text{perm}}\), \(\Phi_{ijkl}=\Phi_{\text{perm}}\),
si puo aggiungere uno strato di sostituzioni simboliche sui tensori in input.

\end{document}
